\documentclass[11pt,a4paper]{article}
\usepackage[margin=24mm]{geometry}
\usepackage{amsmath,amssymb,amsthm,booktabs,hyperref}
\hypersetup{hidelinks,pdftitle={Equality in the Lonely Runner bound for six and seven integer speeds}}
\newtheorem{theorem}{Theorem}
\newtheorem{lemma}[theorem]{Lemma}
\newtheorem{corollary}[theorem]{Corollary}
\newcommand{\LR}{\operatorname{LR}}
\newcommand{\nn}[1]{\lVert #1\rVert}
\setlength{\emergencystretch}{2em}
\begin{document}
\begin{center}
{\LARGE\bfseries Equality in the Lonely Runner bound\\[3pt]for six and seven integer speeds\par}
\medskip
Research draft prepared with OpenAI Codex \quad 26 November 2026
\end{center}
\begin{abstract}
For six distinct positive integer speeds we give a computer-assisted proof that
only dilations of the consecutive speeds attain maximum loneliness $1/7$.
For seven speeds the corresponding equality at $1/8$ has three primitive
solutions. All four examples were known. The issue is completeness at
arbitrarily large speeds. The proof uses lower-dimensional lonely-runner
bounds, a simultaneous approximation, and explicit prime-modulus cover lists.
A short representative of a modular cover either yields a strictly better
time or determines the original speed ratios. Two enumerators are supplied.
\end{abstract}

\section{Equality and the four-replacement question}
Let $V$ consist of $n-1$ distinct positive integers, with a stationary runner
as reference. Write $\nn{x}=\min_{m\in\mathbb Z}|x-m|$,
$[j]=\{1,\ldots,j\}$, and
\[
  \LR(V)=\max_{t\in\mathbb R}\min_{v\in V}\nn{tv}.
\]
An integer dilation preserves $\LR$, since $dt$ traverses the whole circle.
Call a set primitive if its elements have greatest common divisor one.
\begin{theorem}\label{main}
For $n=7,8$, the complete list of primitive $V$ with $\LR(V)\le 1/n$ is
\[
\begin{array}{c|l}
7 &[6]\\
8 &[7],\quad \{1,2,3,4,5,7,12\},\quad \{1,4,5,6,7,11,13\}.
\end{array}
\]
Every listed set attains equality. Every other primitive set admits a time
at which all distances are strictly greater than $1/n$.
\end{theorem}
Goddyn and Wong discuss uniqueness for $n=7$ and list the two exceptional
examples for $n=8$ \cite{GW}. Zhang describes the earlier complete
classifications through five speeds and attributes the example containing
$12$ to Wills \cite[Sec.~1.2]{Z}. The prime-divisor argument below follows
the strategy of Rosenfeld for the ordinary lower bound \cite{R}. For background
on effective bounds and lower-dimensional subtori see \cite{GK,MSS,JA}.

\begin{corollary}\label{four}
Delete exactly four elements of $[n-1]$ and insert four distinct integers
at least $n$. At $n=7$ the only tight result is $3[6]$. At $n=8$ the only
ones are $2[7]$ and $2\{1,2,3,4,5,7,12\}$.
\end{corollary}
Indeed a dilation $dT$ must have exactly $n-5$ entries below $n$.
At least one such entry exists, so $1\le d\le n-1$; direct counting suffices.

\newpage
\section{A strict improvement from a nearby lattice point}
Set $K=n(n-1)/2$. Assume the ordinary lower bound for $n-2$ distinct
positive integer speeds: some time puts each distance at least $1/(n-1)$.
For $n=7$ this is the theorem of Bohman, Holzman and Kleitman \cite{BHK},
and for $n=8$ it is the theorem of Barajas and Serra \cite{BS}.
It applies to fewer speeds by adjoining positive speeds, and to nonzero
rational speeds by clearing denominators, taking absolute values and
removing duplicates.

\begin{lemma}\label{near}
Let $v$ be the positive integer speed vector. If for some $\tau\in\mathbb R$
and $a\in\mathbb Z^{n-1}$,
$\max_i|\tau v_i-a_i|<1/K$, then either $a$ is proportional to $v$
or $\LR(V)>1/n$.
\end{lemma}
\begin{proof}
In the nonproportional case choose indices $i,j$ whose slopes
$a_i/v_i<a_j/v_j$ are adjacent among the distinct slopes. The weighted
average $\lambda=(a_i+a_j)/(v_i+v_j)$ lies strictly between them and equals
no slope. The rational vector $u=a-\lambda v$ has no zero coordinates,
and $u_i=-u_j$. It therefore has at most $n-2$ distinct absolute
coordinates. Choose $s$ so that $\nn{s u_\ell}\ge1/(n-1)$ for every $\ell$.
Put $\theta=s-m$ for a nearest integer $m$ and set
$t=-s\lambda+\theta\tau$. Coordinatewise,
\[
 tv-su=\theta(\tau v-a)-ma .
\]
Modulo integers the displacement has magnitude strictly less than
$1/(2K)$ (and is zero if $\theta=0$). The distance to the nearest integer
is $1$-Lipschitz. Thus every $\nn{tv_\ell}$ exceeds
$1/(n-1)-1/(2K)=1/n$.
\end{proof}

\begin{lemma}\label{bounded}
If $v$ is primitive and $\LR(V)\le1/n$, then
$W=\max_i v_i\le K^{n-2}$.
\end{lemma}
\begin{proof}
Choose a coordinate equal to $W$. For the other $n-2$ coordinates place the
points $(\{jv_i/W\})_{v_i\ne W}$, $0\le j\le K^{n-2}$, in $K^{n-2}$
half-open cubes of side $1/K$ partitioning $[0,1)^{n-2}$. Two lie in the
same cube. Subtracting gives $1\le q\le K^{n-2}$ and integers $a_i$ with
$|qv_i/W-a_i|<1/K$. Use $a_i=q$ at $W$. Lemma~\ref{near} forces
$a=(q/W)v$. Primitivity and B\'ezout's identity give $q/W\in\mathbb Z$.
In particular $W\le q$.
\end{proof}

\begin{lemma}\label{lift}
Suppose $\LR(V)\le1/n$ and there is a unit $h\bmod p$ such that the
absolute least residues of $hv_i\bmod p$, with multiplicity, form a listed
model $T$ with $\max T<p/K$. Then $V$ is an integer dilation of $T$.
\end{lemma}
\begin{proof}
Let $b_i$ be the signed least residues and $a_i=(hv_i-b_i)/p$.
With $\tau=h/p$ their error is $|b_i|/p<1/K$. Lemma~\ref{near} gives
$a=cv$, so $b=(h-pc)v$. The scalar is nonzero. Taking absolute values
identifies $V$ with a positive dilation of $T$. Since $1\in T$, the dilation
factor is itself a speed and is an integer.
\end{proof}

\newpage
\section{The finite prime certificate}
For an odd prime $p$ let $|x|_p$ denote the least absolute residue in
magnitude and let $H=\{1,\ldots,(p-1)/2\}$. Both signs of a speed have
the same bad times. The band of $b\in H$ is
\[
 B(b)=\{j\in H:n|jb|_p\le p\}.
\]
If $\LR(V)\le1/n$, the times $j/p$ imply that the speed bands cover $H$.
Suppose $p$ divides no speed. Multiplying all speeds by the inverse of any
one speed permutes the time classes up to sign. Folding by sign gives a
cover $C\subseteq H$ containing $1$ with at most $n-1$ distinct elements.
Residue collisions are allowed in this reduction.

\medskip\noindent{\bf Finite assertion.}
For each prime in the following row, every cover containing $1$ of size
at most $n-1$ has exactly $n-1$ elements, and the complete list consists of
\begin{equation}\label{list}
 \big\{\,\{|x/b|_p:x\in T\}:T\text{ in the row of
 Theorem~\ref{main}},\ b\in T\,\big\}.
\end{equation}
Division is in $\mathbb F_p$. The list has six covers at each of the
16 primes for $n=7$, and 21 covers at each of the 23 primes for $n=8$.
\begin{center}\small
\begin{tabular}{c p{0.80\linewidth}}\toprule
$n$ & Prime denominators \\ \midrule
7 & 373, 379, 383, 389, 397, 401, 409, 419, 433, 439, 443, 449, 463, 467, 479, 491.\\[3pt]
8 & 367, 379, 383, 389, 397, 419, 431, 463, 487, 499, 503, 509, 521, 541, 547, 557, 563, 569, 587, 593, 599, 601, 607.\\
\bottomrule\end{tabular}
\end{center}
Only these primes are required. For example, the assertion fails for
$n=7$, $p=367$, a case retained as a negative control.

\begin{proof}[Proof of Theorem~\ref{main} from the finite assertion]
Assume $V$ is primitive and $\LR(V)\le1/n$. Lemma~\ref{bounded} gives
$\prod v_i\le(K^{n-2})^{n-1}$. The product of the distinct primes in
each row exceeds this bound. At $n=7$ the bound is below $10^{40}$
and the product exceeds $10^{42}$. At $n=8$ the bound is below $10^{61}$
and the product exceeds $8\cdot10^{61}$. Thus some row prime divides
no speed. The resulting normalized cover is one in \eqref{list}. Its
$n-1$ distinct elements preclude collisions. Undoing the normalization
and multiplying by the listed $b$ gives a unit $h$ as in
Lemma~\ref{lift}. All row primes exceed $K\max T$ (at most 126 and 364
respectively). The lemma and primitivity give $V=T$.

For the converse we check $\LR(T)$ exactly. The lower envelope
$\min_{v\in T}\nn{vt}$ is piecewise linear with nonzero slopes on open
linear pieces. At a maximum a rising and a falling segment meet,
possibly belonging to the same speed. If their speeds are $v,w$, then
$(v+w)t$ is integral. Checking every $t=j/(v+w)$ in $[0,1)$ thus gives
the exact maximum. For the four models in theorem order these are
$1/7,1/8,1/8,1/8$; in each case $t=1/n$ attains the value.
\end{proof}

\newpage
\section{Replaying the enumeration}
Both enumerators construct the bands, fix speed class $1$, and print the
actual sorted speed classes of each cover they find. Neither program is
given the expected templates. A shorter cover would be reported and fail
the comparison. Conversely, if no shorter cover exists, every cover of
the required size is inclusion-minimal and must be enumerated.

\paragraph{Average-coverage search.}
If $G$ times remain and $u$ slots are available, at least one class in any
completion covers at least $G/u$ residual times. The first enumerator
branches on every such class. After a branch, it excludes its class from
later siblings only. Every completion has a first qualifying member in
this loop and follows that branch. At zero residual times the program
emits the partial cover immediately; if no slots remain with uncovered
times it returns. This proves completeness for minimal covers by induction.

\paragraph{Ranked search with an incidence check.}
The second enumerator constructs the time and speed lists using modular
inversion. At each node it ranks the candidates by decreasing number of
residual times hit. A branch chooses the first remaining member of the
completion in this order. If the next $u$ largest counts sum to less than
$G$, every completion in that suffix is impossible. With two slots left,
it chooses one uncovered time and branches on a speed containing it.
The possible partners are the intersection of the speed lists of all
still uncovered times. A class with zero hits cannot be needed by a
minimal cover. The two searches use different final tests and state
representations.

\paragraph{Reproduction and scope.}
Compile either C++17 source in the accompanying directory and pass its
executable to \texttt{verify.py}. The Python driver compares the full rows
with \eqref{list}, independently checks coverage at every nonzero residue,
checks the prime products using integers, and computes the continuous
maxima using fractions. See the repository for commands, both sources and
a replay transcript. This computer-assisted draft awaits external review.
The general equality classification at arbitrary dimension is outside its
scope.

\begin{thebibliography}{9}\small\raggedright
\bibitem{BHK} T. Bohman, R. Holzman and D. Kleitman,
\href{https://doi.org/10.37236/1602}{Six lonely runners}, EJC, 2001.
\bibitem{BS} J. Barajas and O. Serra,
\href{https://arxiv.org/abs/0710.4495}{The lonely runner with seven runners}, EJC, 2008.
\bibitem{GW} L. Goddyn and E. B. Wong,
\href{https://zenodo.org/records/8275490}{Tight instances of the lonely runner}, Integers 6, A38, 2006.
\bibitem{GK} V. Giri and N. Kravitz,
\href{https://doi.org/10.1017/S0305004125101497}{The structure of Lonely Runner spectra}, 2026.
\bibitem{MSS} R. Malikiosis, F. Santos and M. Schymura,
\href{https://doi.org/10.1017/fms.2025.10107}{Linearly-exponential checking is enough for the Lonely Runner Conjecture and some of its variants}, 2025.
\bibitem{R} M. Rosenfeld,
\href{https://arxiv.org/abs/2509.14111}{The lonely runner conjecture holds for eight runners}, 2025.
\bibitem{JA} J. Allikvere,
\href{https://arxiv.org/abs/2609.02604}{Fourteen and fifteen lonely runners}, 2026.
\bibitem{Z} Y. Zhang,
\href{https://arxiv.org/abs/2608.13599}{Tight instances of the Lonely Runner Conjecture: one-entry modifications and the growth of speeds}, v3, 2026.
\end{thebibliography}
\end{document}
